\documentclass[11pt]{article}
\usepackage[a4paper,margin=1in]{geometry}
\usepackage{amsmath,amssymb,amsthm}
\usepackage{graphicx}
\usepackage{microtype}
\usepackage[hidelinks]{hyperref}
\usepackage{enumitem}

\newtheorem{theorem}{Theorem}
\newtheorem{lemma}[theorem]{Lemma}
\newcommand{\B}{\mathcal B}
\newcommand{\Q}{\mathcal Q}
\newcommand{\Kc}{\mathcal K}
\newcommand{\Ad}{\operatorname{Ad}}

\title{A singularizing random-unitary channel whose Cesaro averages diverge\\
\large Counterexample to a mean-ergodic extension suggested in arXiv:2605.24923}
\author{Automated mathematics run \texttt{fa\_banach\_001}, lane 5}
\date{9 August 2026}

\begin{document}
\maketitle

\begin{center}
\fbox{\parbox{0.91\textwidth}{\textbf{Status.}
Candidate counterexample, likely valid, requiring expert human review. It
disproves the universal Cesaro-convergence analogue for singularizing
channels and gives an explicit non-quasi-compact example.}}
\end{center}

\section{Source direction and precise scope}

Dzhenzher studies random-unitary channels defined by Pettis averaging on
possibly nonseparable Hilbert spaces \cite{Dzhenzher}. The conclusion asks for
analogues of strong Cesaro convergence for singularizing channels and points
to quasi-compactness as a possible route.

\begin{figure}[h]
  \centering
  \includegraphics[width=\textwidth]{figures/open_problem_crop.png}
  \caption{The concluding singularizing-channel direction in
  \cite[p.~13]{Dzhenzher}.}
\end{figure}

The source statement is exploratory rather than a numbered yes/no question.
We answer its universal reading:
\begin{quote}
Must the Cesaro averages of a countable-additivity preserving singularizing
random-unitary channel converge in strong operator topology? Is such a channel
necessarily quasi-compact in the measurable-cardinal setting?
\end{quote}
Both answers are negative. The construction uses exactly the source's standing
hypothesis that a cardinal carries an atomless countably additive probability
measure on all of its subsets.

\section{Definitions}

For a Hilbert space $\mathcal H$, a state is a norm-one positive functional on
$\B(\mathcal H)$. A state is \emph{singular} when it vanishes on
$\Kc(\mathcal H)$, the compact operators. A channel $\Q$ is
\emph{singularizing} when $\Q\rho$ is singular for every state $\rho$.

For a group representation $\pi:G\to\mathcal U(\mathcal H)$ and a
countably additive probability $\nu$ on $2^G$, the source's random-unitary
channel is
\[
 \Q_{\nu,\pi}\rho=\int_G \pi(g)\rho\pi(g)^*\,d\nu(g),
\]
where conjugation of a functional is understood by duality. The channel acts
linearly and contractively on $\B(\mathcal H)^*$. Strong operator convergence
of channels therefore means pointwise norm convergence on this dual Banach
space.

\section{Proof intuition}

Use two tensor coordinates for two independent jobs. Atomless averaging of
translations in the first coordinate singularizes every input: each vector
matrix coefficient can see only countably many translations. In the second
coordinate, apply one deterministic bilateral shift at every step. After $n$
iterations the state is supported exactly on the $n$-th second-coordinate
slice, regardless of what happened in the first coordinate. Thus the orbit
states occupy mutually orthogonal corners. Averaging the first $N$ corners
cannot approach averaging the first $2N$: a diagonal sign observable separates
the two means by $1$ for every $N$.

\section{Construction and proof}

Let $K$ be an abelian group of cardinality $\kappa$ with identity $0$.
Assume there is a countably additive probability
\[
 \mu:2^K\longrightarrow[0,1],\qquad \mu(\{j\})=0\quad(j\in K).
\]
This is the Ulam real-valued measurable-cardinal hypothesis used in the
source. Put
\[
 \mathcal H_1=\ell_2(K),\qquad \mathcal H_2=\ell_2(\mathbb Z),\qquad
 \mathcal H=\mathcal H_1\otimes\mathcal H_2.
\]
Write $(e_k)_{k\in K}$ and $(f_m)_{m\in\mathbb Z}$ for the canonical bases.
Let $S_j e_k=e_{j+k}$ and let $Uf_m=f_{m+1}$.

Take $G=K\times\mathbb Z$ and the unitary representation
\[
 \pi(j,m)=S_j\otimes U^m.
\]
Define a probability on all subsets of $G$ by
\[
 \nu(E)=\mu\bigl(\{j\in K:(j,1)\in E\}\bigr).
\]
It is countably additive and is supported on $K\times\{1\}$. The corresponding
channel is
\begin{equation}\label{eq:channel}
 \Q\rho=\int_K (S_j\otimes U)\rho(S_j\otimes U)^*\,d\mu(j).
\end{equation}
By the source's preservation theorem, $\Q$ preserves countable additivity
\cite{Dzhenzher}.

\begin{lemma}\label{lem:singular}
The channel $\Q$ is singularizing.
\end{lemma}

\begin{proof}
Every vector of $\mathcal H\cong\ell_2(K\times\mathbb Z)$ has countable
support. Hence, for fixed $x,v\in\mathcal H$, the set
\[
 E_{x,v}=\{j\in K:\langle v,(S_j\otimes U)x\rangle\neq0\}
\]
is countable: any such $j$ is a difference of first coordinates belonging to
the two countable supports. Therefore $\mu(E_{x,v})=0$.

Let $\rho$ be normal. A positive trace-class operator has a countable spectral
decomposition, so $\rho=\sum_r a_r\rho_{x_r}$ with $a_r\geq0$ and
$\sum_r a_r=1$. For every rank-one projection $P_v$, positivity and monotone
convergence give
\[
 (\Q\rho)(P_v)
 =\int_K\sum_r a_r|\langle v,(S_j\otimes U)x_r\rangle|^2\,d\mu(j)=0,
\]
because the integrand vanishes outside a countable union of countable sets.
Thus $\Q\rho$ vanishes on the compact ideal.

If $\rho$ is singular, then for every compact $A$ and every $j$,
$(S_j\otimes U)^*A(S_j\otimes U)$ is compact, so each integrand in
$(\Q\rho)(A)$ is zero. Finally, every state is a convex combination of a
normal and a singular state by the normal/singular decomposition recalled in
\cite{Dzhenzher}. Hence $\Q\rho$ is singular for every state $\rho$.
\end{proof}

\begin{theorem}\label{thm:main}
The Cesaro averages
\[
 M_N=\frac1N\sum_{n=1}^N\Q^n
\]
do not converge in the strong operator topology on $\B(\mathcal H)^*$.
Indeed, for a vector state $\rho_0$ and every $N\geq1$,
\[
 \|M_N\rho_0-M_{2N}\rho_0\|\geq1.
\]
Moreover, $\Q$ is not quasi-compact.
\end{theorem}

\begin{proof}
Let $\rho_0=\rho_{e_0\otimes f_0}$ and put
$\sigma_n=\Q^n\rho_0$. Every unitary appearing in the $n$-fold integral has
second tensor coordinate $U^n$. Therefore, with
\[
 R_m=I_{\mathcal H_1}\otimes P_{f_m},
\]
we have
\begin{equation}\label{eq:orthogonal}
 \sigma_n(R_m)=\begin{cases}1,&m=n,\\0,&m\neq n.\end{cases}
\end{equation}
For $N\geq1$ define the self-adjoint contraction
\[
 D_N=\sum_{n=1}^N R_n-\sum_{n=N+1}^{2N}R_n.
\]
Equation~\eqref{eq:orthogonal} gives
\[
 (M_N\rho_0)(D_N)=1,
 \qquad
 (M_{2N}\rho_0)(D_N)=0.
\]
Since $\|D_N\|=1$, it follows that
$\|M_N\rho_0-M_{2N}\rho_0\|\geq1$. The vector sequence
$(M_N\rho_0)$ is not Cauchy, so $(M_N)$ cannot converge strongly.

The same projections show that $\|\sigma_n-\sigma_m\|=2$ for $n\neq m$:
test against $R_n-R_m$. Hence the power orbit of $\rho_0$ has no relatively
norm-compact subsequence. A power-bounded quasi-compact operator has relatively
norm-compact power orbits, by its standard finite-dimensional peripheral
spectral decomposition \cite[Uniform ergodic theory]{DS}. Since $\Q$ is a
contraction, it is power bounded, and therefore it cannot be quasi-compact.
\end{proof}

\section{Verification, scope, and novelty}

The construction remains entirely inside the source's class: $\nu$ is a
countably additive probability on a group, $\pi$ is unitary, and
Equation~\eqref{eq:channel} is its Pettis random-unitary channel. The
large-cardinal measure is not an extra unproved dependency; it is the stated
hypothesis under which the source constructs its singularizing channels.

The script \texttt{code/check\_cesaro\_gap.py} verifies with exact rational
arithmetic, for $N=1,\ldots,100$, the scalar weight calculation behind the
diagonal witness. It returns gap $1$ in every case. This is a finite sanity
check, not a substitute for the proof.

A bounded novelty search on 9 August 2026 covered the run registry and indexed
web/arXiv searches for the source id and title, singularizing quantum channels,
Cesaro nonconvergence, mean ergodicity, and quasi-compactness. It found the
source paper and general finite-dimensional quantum ergodicity literature,
but no prior result matching this construction or resolving the source's
singularizing direction. Novelty confidence is moderate to high, but the
packet should remain a candidate counterexample until expert review.

The result rules out an unrestricted analogue. It does not preclude strong
Cesaro convergence or quasi-compactness under additional tightness,
compactness, or mixing assumptions.

\begin{thebibliography}{9}

\bibitem{Dzhenzher}
S. V. Dzhenzher,
\emph{Singularizing preserving countable additivity quantum channels on
quantum measurable cardinals}, arXiv:2605.24923 (2026).

\bibitem{DS}
N. Dunford and J. T. Schwartz,
\emph{Linear Operators, Part I: General Theory},
Wiley-Interscience, 1958.

\end{thebibliography}

\end{document}
